% Part: counterfactuals
% Chapter: introduction
% Section: counterfactuals

\documentclass[../../../include/open-logic-section]{subfiles}

\begin{document}

\olfileid{cnt}{int}{cnt}

\olsection{الشرطيات المخالفة للواقع}

يُبنى شكل شائع ومهم جدًا من تراكيب «إذا \dots، فإن \dots» في الإنجليزية
باستعمال صيغة الماضي الالتزامية من فعل \emph{يكون}: «لو كان الأمر أن
\dots، لكان الأمر أن \dots». ولأن مقدّم مثل هذا الشرط يكون كاذبًا عادة،
أي مخالفًا للواقع، تسمى هذه التراكيب \emph{شرطيات مخالفة للواقع}
(ولأنها تستعمل الصيغة الالتزامية من فعل \emph{يكون}، تسمى أيضًا
\emph{شرطيات التزامية}). وتتميز من الشرطيات \emph{الخبرية} التي تأتي
على صورة «إذا كان الأمر أن \dots، فإن الأمر أن \dots». وتختلف الشرطيات
المخالفة للواقع عن الشرطيات الخبرية في شروط صدقها. انظر إلى مثال آدامز
الشهير:
\begin{quote}
  إذا لم يقتل أوزوالد كينيدي، فقد قتله شخص آخر.
  
  لو لم يكن أوزوالد قد قتل كينيدي، لكان شخص آخر قد قتله.
\end{quote}
الأولى خبرية، والثانية مخالفة للواقع. والأولى صادقة بوضوح: فنحن نعلم أن
الرئيس جون ف. كينيدي قتله \emph{شخص ما}؛ وإذا لم يكن ذلك الشخص
(خلافًا لتقرير وارن) لي هارفي أوزوالد، فقد قتل كينيدي شخص آخر. أما
الثانية فتقول شيئًا مختلفًا. فهي تدعي أنه لو لم يكن أوزوالد قد قتل
كينيدي، أي لو تم تجنب إطلاق النار في دالاس أو لم ينجح، لكان التاريخ قد
جرى بعد ذلك على نحو تنجح فيه محاولة اغتيال أخرى. ولكي تصدق، فلا بد أن
تكون قوى نافذة قد تآمرت لضمان موت جون ف. كينيدي (كما يعتقد كثير من
أصحاب نظريات المؤامرة حول اغتياله).

لا يزال ثمة نقاش قائم حول ما إذا كان الشرط \emph{الخبري} يُصاغ على نحو
صحيح بالشرط المادي، وبخاصة حول إمكان «تفسير» مفارقات الشرط المادي
بطريقة تنسجم مع كونه يحدد شروط صدق الشرطيات الخبرية في الإنجليزية.
أما أن الشرطيات المخالفة للواقع لا يمكن تمثيلها تمثيلًا صحيحًا بالشروط
المادية فليس موضع خلاف. وهذا واضح لأن مقدمات الشرطيات المخالفة للواقع،
وإن كانت كاذبة في العموم، لا تجعل كل شرط مخالف للواقع ذا مقدّم كاذب
صادقًا؛ فشرط آدامز المخالف للواقع مثال على ذلك إذا صدقنا تقرير وارن
ولم تكن هناك مؤامرة لاغتيال كينيدي.

تؤدي الشرطيات المخالفة للواقع دورًا مهمًا في الاستدلال السببي؛ فمن أبرز
استعمالاتها التعبير عن العلاقات السببية. فمثلًا، يسبب حك عود الثقاب
اشتعاله، ويمكن التعبير عن ذلك بالقول: «لو حُك هذا العود، لاشتعل».
ولا يمكن استعمال الشرطيات المادية، ولا الشرطيات الخبرية في العموم،
للتعبير عن ذلك: فالعبارة «يُحك عود الثقاب~$\lif$ يشتعل عود الثقاب»
تصدق إذا لم يُحك العود قط، بصرف النظر عما كان سيحدث لو حُك. والأسوأ أن
العبارة «يُحك عود الثقاب~$\lif$ يتحول عود الثقاب إلى باقة زهور» تصدق
هي أيضًا إذا لم يُحك قط، مع أنه يقينًا لن يتحول إلى باقة زهور لو حُك.

ولا يزال المنطق الصحيح للشرطيات المخالفة للواقع موضع نقاش. وقد قدّم
ستالنيكر ولويس تحليلًا مؤثرًا لها. ووفقًا لهما، يصدق الشرط المخالف
للواقع «لو كان الأمر أن~$S$، لكان الأمر أن~$T$» إذا وفقط إذا كانت~$T$
صادقة في الحالة المخالفة للواقع («العالم الممكن») الأقرب إلى حال العالم
الفعلي والتي تكون فيها~$S$ صادقة. ويسمى هذا تحليلًا «وجوديًا» لأنه
يرجع إلى أنطولوجيا للعوالم الممكنة. وتستعمل تحليلات أخرى الاحتمالات
الشرطية أو نظريات مراجعة الاعتقاد. وقد تكاثرت المنطقيات المختلفة
المقترحة للشرطيات المخالفة للواقع. بل لا يوجد منطق واحد للشرطيات
المخالفة للواقع عند لويس وستالنيكر: فرغم أنهما اقترحا تحليلين متقاربين
يرجعان إلى أقرب العوالم الممكنة، فإن افتراضاتهما تفضي إلى استدلالات
صالحة مختلفة.

\end{document}
